\newcommand{\outofk}{{$3$-out-of-$k\!+\!1$}\xspace}

In this section, we formally define the relaxed computational problem of Sperner's lemma, and its variants that we will use for the applications on envy-free cake cutting.  
In particular, we state our main technical result \cref{thm:main}, deferring the proof to \cref{sec:warmup-proof-for-approx-sperner,,sec:final-proof-for-approx-sperner}.

Here, we use the term {\em approximate Sperner's lemma} for the following relaxed version of the original Sperner's lemma~\cite{sperner1928neuer}: in any triangulation of a large $k$-dimensional simplex with a $k+1$-Sperner coloring that satisfies appropriate boundary condition, there always exists a small triangle (a.k.a., $2$-simplex) whose three vertices are colored differently (because Sperner's lemma states that there exists a small $k$-simplex whose $k+1$ vertices are colored by all $k+1$ colors).
Formally, we define the following \outofk Approximate \Sperner problem. 


\begin{definition}[\outofk Approximate \Sperner]\label{def:approx-kdsperner}
    In \outofk Approximate \Sperner (with a triangulation side-length of $2^{-n}$), we are given a function 
    $C:\Delta^k_n \to [k+1]$ satisfying $C(\vec{x}) \in \mathcal{I}_{>0}(\vec{x})$ for any $\vec{x}\in \Delta^k_n$.
    The goal of \outofk Approximate \Sperner is to find three points $\vec{x_1},\vec{x_2}, \vec{x_3}\in \Delta_n^k$ such that
    \begin{enumerate}[itemsep=0.2em,topsep=0.5em]
        \item {\bf the three points are close to each other}, i.e., $\forall i,j\in [3], \normtxt{\vec{x_i}-\vec{x_j}}{\infty}\leq 2^{-n}$; and 
        \item {\bf the induced triangle is trichromatic}, i.e., $|\{C(\vec{x_i}): i\in [3]\}|=3$.
    \end{enumerate}
\end{definition} 

\begin{remark}
    \label{remark:equiv-at-2}
    When $k=2$, the \outofk Approximate \Sperner problem is simply equivalent to the $k${\rm D}-\Sperner problem.
\end{remark}
    
To apply the hardness result of \outofk Approximate \Sperner to the envy-free cake-cutting problem, we further consider the special case 
where the coloring defined by $C$ satisfies the following symmetry constraints on the boundary of the simplex. 
First, we define the symmetry that naturally arises from cake-cutting valuations.
For example, suppose we cut the cakes using two cuts in two different ways: with cut locations $(0.1, 0.1)$ and $(0.1, 1)$. Both partitions of the cake have the same two non-trivial pieces: $[0,0.1]$ and $[0.1,1]$. Thus we expect that the preference for each agent between those pieces is the same regardless of the representation on the simplex. 
This inspires us to define a symmetry constraint on cuts that have the same set of non-zero pieces.


\begin{definition}
    \label{def:indexing}
    Indexing $\idx(\vec{a}, v)$ of $v$ in an array $\vec{a}$ is defined as the index $i$ such that $a_i=v$. 
\end{definition}

\begin{restatable}[symmetric points in a coloring]{definition}{symmetrydef}
    \label{def:symmetry-in-coloring}
    We say $\vec{x}$ and $\vec{y}$ are symmetric in a coloring $C$ if 
    \begin{itemize}[topsep=0.5em, itemsep=0.1em]
        \item the number of non-trivial entries in $\vec{x}$ and $\vec{y}$ are the same, i.e., $|\mathcal{I}_{>0}(\vec{x})| = |\mathcal{I}_{>0}(\vec{y})|$.
        \item the indexing of $C$ are the same for the arrays of non-trivial indices of $\vec{x}$ and $\vec{y}$, i.e., $\idx(\mathcal{I}_{>0}(\vec{x}), C(\vec{x})) = \idx(\mathcal{I}_{>0}(\vec{y}), C(\vec{y}))$.
    \end{itemize}
    We use $\vec{x}\sim_C \vec{y}$ to denote this symmetry. 
\end{restatable}

Because the two conditions for symmetry are both equations, which satisfy reflexivity, symmetry, and transitivity, this symmetry is a valid equivalence relation. 
\begin{fact}
    \label{lem:symmetry-is-equiv-relation}
    For any coloring $C$, the symmetry of points in coloring $C$ is an equivalence relation, i.e., a binary relation satisfying reflexive ($\vec{x}\sim_{C} \vec{x}$), symmetric ($\vec{x}\sim_{C} \vec{y} \rightarrow \vec{y}\sim_{C} \vec{x}$) and transitive ($\vec{x}\sim_{C} \vec{y}$ and $\vec{y}\sim_{C} \vec{z} \rightarrow \vec{x}\sim_{C} \vec{z}$).
\end{fact}

In the symmetric version of the \outofk Approximate \Sperner problem, the input guarantees that symmetric points on the boundary are symmetric in the input coloring. 
We say two points on the boundary are symmetric to each other, if the resulting vectors are the same after we remove the zero entries in them.
\begin{restatable}[\outofk Approximate Symmetric \Sperner]{definition}{symsperner}
    \label{def:approx-symmetric-kd-sperner}
    \outofk Approximate Symmetric \Sperner is a special case of \outofk Approximate \Sperner, where the given circuit $C$ further satisfies the following property: for any $\vec{x}\in \Delta_n^{k}$ and any $i,j \in [k+1]$, $(\vec{x}_{1:i-1}, 0, \vec{x}_{i:k}) \sim_C (\vec{x}_{1:j-1},0, \vec{x}_{j:k})$. 
\end{restatable}

Our main technical result is the \PPAD-hardness of this \outofk Approximate Symmetric \Sperner problem. 
Because this problem is clearly a member of the class \PPAD, its complexity is \PPAD-complete. 

\begin{restatable}[]{theorem}{mainthm}
    \label{thm:main}
    For and any $k=\poly(n)$, \outofk Approximate Symmetric \Sperner
    with a triangulation side-length of $2^{-4kn}$ 
    is \PPAD-complete.
    Further, for any $k\geq 2$, it requires a query complexity of $2^{\Omega(n)} / \poly(n,k)$.
\end{restatable}